% Hodina 52, 1. ročník SPŠ. Rozsah: Rozpis_vyuky_MAT_1_rocnik.tex, řádek 52.

% Plán: 1_rocnik/tematicke_plany/SPS.pdf; společná projektová šablona.

% Učebnice: ucebnice/01_Zakladni_poznatky/Algebraicke_vyrazy/Algebraické výrazy.tex

% Původ úloh: P1: učebnice, „Úpravy lomených výrazů“, příklad x^2-4 a x^2-2x. P2: stejný výraz z příkladu „Definiční obor“, doplněna vlastní úprava. Ostatní vlastní.

% Metodický zdroj: externi_zdroje/realisticky/Lomene_vyrazy/01_Kraceni_a_rozsirovani_lomenych_vyrazu.pdf

% URL: http://www.realisticky.cz/ucebnice/01%20Matematika%20S%C5%A0/01%20Z%C3%A1kladn%C3%AD%20poznatky/09%20Lomen%C3%A9%20v%C3%BDrazy/01%20Kr%C3%A1cen%C3%AD%20a%20roz%C5%A1i%C5%99ov%C3%A1n%C3%AD%20lomen%C3%BDch%20v%C3%BDraz%C5%AF.pdf

\documentclass[11pt,a4paper]{article}

% Spolecny styl pro ucitelske pripravy MAT (Work/MAT)
% Verze 1.1, 2026-10-23
\NeedsTeXFormat{LaTeX2e}
\ProvidesPackage{priprava}[2026/10/23 v1.1 Sablona ucitelskych priprav MAT]

% Zakladni typografie a cestina
\RequirePackage[T1]{fontenc}
\RequirePackage[utf8]{inputenc}
\RequirePackage[czech]{babel}
\RequirePackage{lmodern}
\RequirePackage{microtype}
\RequirePackage{amsmath,amssymb}
\RequirePackage{xcolor}
\RequirePackage{enumitem}
\RequirePackage[most]{tcolorbox}
\RequirePackage{tabularx,booktabs,array}
\RequirePackage{fancyhdr}
\RequirePackage{needspace}
\RequirePackage{geometry}
\RequirePackage[hidelinks]{hyperref}

\geometry{a4paper,top=17mm,bottom=18mm,left=23mm,right=23mm}
\definecolor{notegray}{gray}{0.39}
\setlength{\parindent}{0pt}
\setlength{\parskip}{2.5pt plus 1pt minus 0.5pt}
\setlist[itemize]{topsep=2pt,itemsep=1pt,parsep=0pt,leftmargin=1.45em}
\setlist[enumerate]{topsep=2pt,itemsep=2pt,parsep=0pt,leftmargin=1.65em}
\setcounter{secnumdepth}{0}
\renewcommand{\arraystretch}{1.12}

% Zahlavi: vlevo rocnik | globalni tema, vpravo cislo: nazev hodiny
% Zapati zustava zcela prazdne.
\pagestyle{fancy}
\fancyhf{}
\renewcommand{\headrulewidth}{0.35pt}
\renewcommand{\footrulewidth}{0pt}
\setlength{\headheight}{14pt}
\newcommand{\nastavPripravu}[4]{%
  \fancyhead[L]{\small #1\enspace\textbar\enspace #2}%
  \fancyhead[R]{\small #3: #4}%
}

% Minimalisticke cernobile boxy. Kazda definice zvlast.
\tcbset{commonlesson/.style={enhanced,breakable=false,colback=white,colframe=black,
  colbacktitle=white,coltitle=black,boxrule=0.45pt,
  arc=0.8mm,outer arc=0.8mm,left=3mm,right=3mm,top=1.6mm,bottom=1.6mm,
  toptitle=1mm,bottomtitle=0.7mm,fonttitle=\bfseries,
  attach boxed title to top left={xshift=2.5mm,yshift=-1.15mm},
  boxed title style={colback=white,colframe=white,boxrule=0pt,left=1mm,right=1mm}}}
\newtcolorbox{zadaniBox}[1][]{commonlesson,title={Zadání},#1}
\newtcolorbox{prikladBox}[1][]{commonlesson,title={Příklad},#1}
\newtcolorbox{definiceBox}[1][]{commonlesson,title={Definice},#1}

% Spolecny box pro zadani a jeho reseni, bez mezery mezi castmi.
% Pouziti:
%   \begin{prikladReseniBox}
%     Zadani ulohy.
%     \tcblower
%     \begin{align*} ... \end{align*}
%   \end{prikladReseniBox}
% Spodni cast ma jemny sedy podklad a delici linku; lze v ni sazet
% vice radku matematiky. Obe casti maji jen jeden spolecny ramecek.
\tcbset{solutionlesson/.style={
  commonlesson,
  enhanced jigsaw,
  breakable,
  colbacklower=black!3,
  segmentation style={solid,draw=black!25,line width=0.35pt},
  fontlower=\small\color{notegray},
  before lower={\noindent\textbf{Řešení:}\enspace},
  leftlower=3mm,rightlower=3mm
}}
\newtcolorbox{prikladReseniBox}[1][]{solutionlesson,title={Příklad},#1}
\newtcolorbox{zadaniReseniBox}[1][]{solutionlesson,title={Zadání},#1}

% Pomucky pro autora pripravy: sedive poznamky; \reseni pro starsi pripravy
\newcommand{\poznamka}[1]{\par\begingroup\small\color{notegray}\itshape
  \textnormal{\textbf{Poznámka pro učitele.}}\enspace #1\par\endgroup}
\newcommand{\reseni}[1]{\par\begingroup\small\color{notegray}
  \textbf{Očekávané řešení:}\enspace #1\par\endgroup}
\newcommand{\blok}[1]{\needspace{5\baselineskip}\subsection*{#1}}
\newcommand{\vyrok}[1]{\uv{#1}}


\nastavPripravu{1. ročník}{Algebraické výrazy}{52}{Lomené výrazy: podmínky a krácení}

\begin{document}

\blok{Výklad a souvislosti}

Zlomek s proměnnou ve jmenovateli může být pro některá dosazení nedefinovaný. Podmínky proto stanovíme dříve, než výraz rozložíme a zkrátíme. Rozklad na součin slouží k nalezení celých společných činitelů.

\begin{definiceBox}[title={Definice},breakable=false]
\textbf{Racionální lomený výraz} je podíl mnohočlenů $A(x)/B(x)$, kde $B$ není nulový mnohočlen. Jeho hodnoty jsou definovány právě pro $B(x)\ne0$.
\end{definiceBox}


\begin{definiceBox}[title={Definice},breakable=false]
\textbf{Krácení lomeného výrazu} je odstranění společného nenulového činitele čitatele a jmenovatele: $AC/(BC)=A/B$ při $B\ne0$, $C\ne0$.
\end{definiceBox}


Krátíme činitele, nikoli jednotlivé sčítance. Rozšiřování je opačný postup: násobíme čitatele i jmenovatele stejným nenulovým výrazem. Po krácení ponecháme původní podmínky, i když z výsledného jmenovatele už některý zákaz nevyplývá.

\blok{Řešené příklady}

\begin{prikladReseniBox}[title={Příklad 1},breakable=false]
Zkrať $\frac{x^2-4}{x^2-2x}$ a urči podmínky.
\tcblower
Jmenovatel je $x(x-2)$, proto $x\ne0,2$. Na tomto oboru $\frac{(x-2)(x+2)}{x(x-2)}=\frac{x+2}{x}$. Hodnota $x=2$ zůstává zakázána.
\end{prikladReseniBox}


\begin{prikladReseniBox}[title={Příklad 2},breakable=false]
Uprav $\frac{x^2-9}{x^2-5x+6}$ a určuj obor před krácením.
\tcblower
Rozklady jsou $(x-3)(x+3)$ a $(x-2)(x-3)$. Proto $x\ne2,3$ a výsledek je $\frac{x+3}{x-2}$ při obou původních zákazech.
\end{prikladReseniBox}


\blok{Samostatné procvičení}

\begin{zadaniReseniBox}[title={Procvičení},breakable=false]
\begin{enumerate}[label=\alph*)]
\item Zkrať $\frac{x^2-1}{x^2+x}$.
\item Zkrať $\frac{2x+2}{x+1}$.
\item Vysvětli, proč nelze „vykrátit $x$“ ve výrazu $\frac{x+2}{x+3}$.
\end{enumerate}
\tcblower
\begin{enumerate}[label=\alph*)]
\item $\frac{(x-1)(x+1)}{x(x+1)}=\frac{x-1}{x}$ při $x\ne0,-1$.
\item $2$ při $x\ne-1$; společným činitelem je celá závorka $x+1$.
\item $x$ není společným činitelem celých součtů. Chybné odstranění $x$ by dalo $2/3$, ale například při $x=1$ je původní hodnota $3/4$. Původní obor vyžaduje $x\ne-3$.
\end{enumerate}
\end{zadaniReseniBox}


\poznamka{Ke každému výrazu piš podmínky na stejný řádek jako výsledek. Nepřeváděj je jen na podmínky zjednodušeného tvaru. V dalším procvičování vrať vždy alespoň jednu zakázanou hodnotu, která po krácení zmizela ze jmenovatele.}

\end{document}
